% =============================================================================
% SYSTEM DESIGN AND ARCHITECTURE
% =============================================================================
% This section is self-contained and can be pasted into the main thesis document.
%
% REQUIRED PACKAGES (add to preamble):
%   \usepackage{tikz}
%   \usetikzlibrary{shapes.geometric, arrows.meta, positioning, fit, backgrounds, calc}
%   \usepackage{booktabs}
%   \usepackage{hyperref}
%
% APPENDIX CONTENTS (referenced but not included here):
%   - Appendix A: Architecture Decision Records (ADRs) - full text
%   - Appendix B: API Specifications and JSON-RPC schemas
%   - Appendix C: Scout Catalog JSON schema
%
% Implementation details are available at: [GitHub Repository URL]
% =============================================================================

% TikZ style definitions for architecture diagrams
\tikzset{
    % Layer boxes
    layer/.style={
        rectangle,
        rounded corners=3pt,
        draw=#1,
        fill=#1!10,
        line width=1pt,
        minimum width=12cm,
        minimum height=2cm,
        align=center,
        font=\small
    },
    % Component boxes within layers
    component/.style={
        rectangle,
        rounded corners=2pt,
        draw=#1,
        fill=#1!5,
        line width=0.5pt,
        minimum height=0.6cm,
        align=center,
        font=\footnotesize
    },
    % Pipeline stages
    stage/.style={
        rectangle,
        rounded corners=3pt,
        draw=black!60,
        fill=blue!8,
        line width=0.8pt,
        minimum width=4cm,
        minimum height=0.9cm,
        align=center,
        font=\small
    },
    % Stage output annotation
    stageout/.style={
        font=\footnotesize\itshape,
        text=black!70,
        anchor=west
    },
    % Network node
    netnode/.style={
        rectangle,
        rounded corners=4pt,
        draw=#1,
        fill=#1!12,
        line width=1pt,
        minimum width=7cm,
        minimum height=2.2cm,
        align=center,
        font=\small
    },
    % Arrow styles
    dataarrow/.style={
        -{Stealth[length=2.5mm]},
        line width=0.8pt,
        draw=black!70
    },
    protocollabel/.style={
        font=\footnotesize,
        fill=white,
        inner sep=2pt
    }
}

\section{System Design and Architecture}
\label{sec:system-architecture}

This section presents the architectural design of the natural language to SQL system.
The architecture addresses the core challenges of enabling non-technical users to query enterprise databases autonomously, while maintaining security, traceability, and deployability in small and medium enterprise (SME) environments.
We first establish the design principles that guided architectural decisions, then present the high-level system structure, followed by detailed descriptions of the key subsystems: autonomous schema modeling, the query reasoning pipeline, and production integration mechanisms.

% -----------------------------------------------------------------------------
\subsection{Design Principles}
\label{subsec:design-principles}
% -----------------------------------------------------------------------------

The system architecture is guided by five core principles that emerged from the requirements analysis and iterative development process:

\begin{description}
    \item[Autonomy.] The system must operate without requiring users to possess database knowledge. This extends beyond query translation to include autonomous discovery of database schemas, understanding of table relationships, and generation of appropriate join strategies. Users should be able to ask questions in natural language without understanding the underlying data model.

    \item[Schema-Agnosticism.] Rather than being hardcoded to a specific database schema, the system must dynamically adapt to arbitrary enterprise databases. This is achieved through a semantic catalog that is built at startup and continuously refined, enabling the system to reason about previously unseen tables and columns.

    \item[Read-Only Safety.] Given that the system operates on production databases containing business-critical data, all database interactions are constrained to read-only operations. This constraint is enforced at multiple architectural layers: the query validation layer, the database connection configuration, and the database user permissions.

    \item[Traceability.] Every step of the query processing pipeline is logged and can be audited. This includes the original natural language query, extracted intent, discovered tables, generated SQL, and execution results. Traceability serves both debugging purposes during development and compliance requirements in production deployments.

    \item[SME Deployability.] The architecture must be deployable in typical SME infrastructure, which often includes legacy systems, VPN-based network access, and limited IT resources. This influenced decisions around proxy-based database access and the separation of components across network boundaries.
\end{description}

These principles are operationalized through the architectural decisions documented in the project's Architecture Decision Records (ADRs), which provide detailed rationale for each design choice.\footnote{The complete set of ADRs is available in the project repository and summarized in Appendix~\ref{appendix:adrs}.}

% -----------------------------------------------------------------------------
\subsection{High-Level Architecture}
\label{subsec:high-level-architecture}
% -----------------------------------------------------------------------------

The system follows a layered architecture with clear separation of concerns between user interaction, reasoning and orchestration, and data access. Figure~\ref{fig:high-level-architecture} illustrates the three primary layers and their interactions.

\begin{figure}[htbp]
    \centering
    \begin{tikzpicture}[node distance=0.8cm]
        % Presentation Layer
        \node[layer=blue!70] (presentation) {
            \textbf{Presentation Layer}\\[2pt]
            \begin{tikzpicture}
                \node[component=blue!70] (ui) {Web UI};
                \node[component=blue!70, right=0.3cm of ui] (rest) {REST API};
                \node[component=blue!70, right=0.3cm of rest] (ws) {WebSocket Streaming};
            \end{tikzpicture}
        };

        % Protocol label 1
        \node[below=0.15cm of presentation, protocollabel] (proto1) {HTTP / WebSocket};

        % Reasoning Layer
        \node[layer=orange!70, below=0.7cm of presentation] (reasoning) {
            \textbf{Reasoning \& Orchestration Layer}\\[2pt]
            \begin{tikzpicture}
                \node[component=orange!70] (agent) {ReAct Agent};
                \node[component=orange!70, right=0.2cm of agent] (intent) {Intent Parser};
                \node[component=orange!70, right=0.2cm of intent] (ranker) {Table Ranker};
                \node[component=orange!70, right=0.2cm of ranker] (sqlgen) {SQL Generator};
                \node[component=orange!70, right=0.2cm of sqlgen] (logger) {Debug Logger};
            \end{tikzpicture}
        };

        % Protocol label 2
        \node[below=0.15cm of reasoning, protocollabel] (proto2) {JSON-RPC 2.0 (MCP)};

        % Data Access Layer
        \node[layer=green!60, below=0.7cm of reasoning] (dataaccess) {
            \textbf{Data Access Layer (MCP Server)}\\[2pt]
            \begin{tikzpicture}
                \node[component=green!60] (discovery) {Discovery Tools};
                \node[component=green!60, right=0.2cm of discovery] (executor) {Query Executor};
                \node[component=green!60, right=0.2cm of executor] (guard) {Safety Guardrails};
                \node[component=green!60, right=0.2cm of guard] (scout) {Scout Runner};
                \node[component=green!60, right=0.2cm of scout] (catalog) {Catalog Store};
            \end{tikzpicture}
        };

        % Protocol label 3
        \node[below=0.15cm of dataaccess, protocollabel] (proto3) {ODBC / Native Driver};

        % Database Layer
        \node[layer=red!50, below=0.7cm of dataaccess] (database) {
            \textbf{Enterprise Database}\\[2pt]
            \begin{tikzpicture}
                \node[component=red!50] (sqlserver) {SQL Server};
                \node[component=red!50, right=0.5cm of sqlserver] (postgres) {PostgreSQL};
                \node[component=red!50, right=0.5cm of postgres] (erp) {ERP Data (943 Tables)};
            \end{tikzpicture}
        };

        % Arrows between layers
        \draw[dataarrow] (presentation.south) -- (reasoning.north);
        \draw[dataarrow] (reasoning.south) -- (dataaccess.north);
        \draw[dataarrow] (dataaccess.south) -- (database.north);

    \end{tikzpicture}
    \caption{High-level system architecture showing the four primary layers: presentation, reasoning and orchestration, data access (MCP server), and enterprise database. Communication between the orchestration layer and data access layer uses the Model Context Protocol (MCP) over JSON-RPC 2.0.}
    \label{fig:high-level-architecture}
\end{figure}

\subsubsection{Presentation Layer}

The presentation layer provides the user interface through which natural language queries are submitted and results are displayed. The primary interface is a web-based chat application that supports real-time streaming of intermediate reasoning steps, allowing users to observe the system's decision-making process. The interface is intentionally minimal, requiring only a text input field and a response area, reinforcing the principle that users need not understand database concepts to use the system.

\subsubsection{Reasoning and Orchestration Layer}

The reasoning layer implements the core intelligence of the system using a ReAct (Reasoning and Acting) agent pattern~\cite{yao2023react}. Built on the LangGraph framework, this layer orchestrates the query processing pipeline through iterative cycles of reasoning about the current state, selecting and executing tools, and observing results. The agent has access to a bounded set of tools for database discovery and query execution, described in Section~\ref{subsec:query-pipeline}.

A key architectural decision was to implement a \emph{simple agent} architecture rather than a multi-agent orchestration system. Early iterations explored specialized agents for intent parsing, table discovery, join planning, and SQL generation. However, empirical evaluation revealed that a single ReAct agent with composite tools achieved comparable accuracy with significantly reduced latency and system complexity. This decision is documented in ADR-0030 and ADR-0031.

\subsubsection{Data Access Layer}

The data access layer implements the Model Context Protocol (MCP), providing a standardized interface for all database interactions. This layer runs as an independent service, potentially on a separate machine with direct database connectivity. The separation enables deployment scenarios where the reasoning layer runs on a developer workstation while the data access layer runs on a server within the enterprise network.

The MCP server exposes database functionality through a JSON-RPC 2.0 interface with the following core operations:

\begin{itemize}
    \item \texttt{search\_tables(query, limit)} -- Semantic search over the schema catalog
    \item \texttt{describe\_table(table\_name)} -- Detailed table metadata including columns and relationships
    \item \texttt{get\_column\_index(tables)} -- Exact column names for specified tables
    \item \texttt{list\_relations(table\_name)} -- Foreign key relationships for join planning
    \item \texttt{execute\_query(sql, limit)} -- Bounded SQL execution with safety validation
\end{itemize}

All operations enforce read-only constraints and implement configurable limits on result set sizes (default: 1000 rows) and execution time (default: 60 seconds).

% -----------------------------------------------------------------------------
\subsection{Autonomous Schema Modeling (Scout Mode)}
\label{subsec:scout-mode}
% -----------------------------------------------------------------------------

A distinguishing feature of the system is its ability to autonomously model enterprise database schemas without manual configuration. This capability, termed \emph{Scout Mode}, addresses the schema-agnosticism principle by building a semantic catalog that enables intelligent table discovery and query planning.

\subsubsection{Motivation}

Enterprise databases, particularly those underlying ERP systems, often contain hundreds or thousands of tables with domain-specific naming conventions. The database used for evaluation in this work contains 943 tables across multiple schemas, with German-language table names following vendor-specific prefixes (e.g., \texttt{KHKArtikel}, \texttt{KHKLagerplatzbestaende}). Without schema modeling, each user query would require exhaustive database introspection, introducing latencies of 10--50 seconds per query---unacceptable for interactive use.

\subsubsection{Catalog Construction}

Scout Mode executes asynchronously at MCP server startup, performing the following steps:

\begin{enumerate}
    \item \textbf{Schema Introspection.} Query the database's \texttt{information\_schema} to enumerate all tables, views, columns, data types, and foreign key relationships.

    \item \textbf{Metadata Enrichment.} For each table, compute derived metadata including:
    \begin{itemize}
        \item Estimated row count (from database statistics)
        \item Column classification by type (numeric, date, text)
        \item Foreign key connectivity (count of inbound/outbound relationships)
        \item Primary key identification
    \end{itemize}

    \item \textbf{Semantic Description Generation.} Automatically generate human-readable descriptions for tables based on structural analysis. For example, a table with columns named \texttt{Artikelnummer}, \texttt{Bezeichnung}, and \texttt{Preis} is described as ``Product master data with article identifiers and pricing information.''

    \item \textbf{Index Construction.} Build in-memory indexes for efficient fuzzy matching, including normalized table names with vendor prefixes stripped (e.g., \texttt{KHKArtikel} $\rightarrow$ \texttt{Artikel}).

    \item \textbf{Persistence.} Serialize the catalog to disk with GZIP compression for fast subsequent startups.
\end{enumerate}

The catalog construction process completes in 30--120 seconds depending on database size, but this cost is amortized across a configurable time-to-live (TTL) period of 7 days, after which the catalog is automatically refreshed.

\subsubsection{Catalog Structure}

The semantic catalog organizes metadata hierarchically, as illustrated in Table~\ref{tab:catalog-structure}. The complete JSON schema is provided in Appendix~\ref{appendix:catalog-schema}.

\begin{table}[htbp]
    \centering
    \caption{Scout catalog structure showing metadata captured for each table.}
    \label{tab:catalog-structure}
    \begin{tabular}{lp{8cm}}
        \toprule
        \textbf{Field} & \textbf{Description} \\
        \midrule
        \texttt{full\_name} & Fully qualified table name (e.g., \texttt{dbo.KHKArtikel}) \\
        \texttt{type} & Object type (\texttt{TABLE} or \texttt{VIEW}) \\
        \texttt{estimated\_rows} & Row count estimate from database statistics \\
        \texttt{columns} & List of column definitions with types and constraints \\
        \texttt{numeric\_columns} & Subset of columns suitable for aggregation \\
        \texttt{date\_columns} & Subset of columns suitable for temporal filtering \\
        \texttt{fk\_count} & Number of foreign key relationships \\
        \texttt{semantic\_description} & Auto-generated natural language description \\
        \texttt{role\_hints} & Inferred column roles (e.g., \texttt{id}, \texttt{amount}, \texttt{status}) \\
        \bottomrule
    \end{tabular}
\end{table}

\subsubsection{Semantic Search and Ranking}

When the reasoning layer requests table discovery, the Scout catalog enables semantic search rather than simple string matching. The ranking algorithm combines multiple signals:

\begin{equation}
    \text{score}(t, q) = w_1 \cdot \text{exact}(t, q) + w_2 \cdot \text{fuzzy}(t, q) + w_3 \cdot \text{column}(t, q) + w_4 \cdot \text{intent}(t, q) + \text{penalties}(t)
    \label{eq:ranking}
\end{equation}

\noindent where $t$ is a table, $q$ is the search query, and:

\begin{itemize}
    \item $\text{exact}(t, q)$ returns 1.0 for exact substring matches
    \item $\text{fuzzy}(t, q)$ computes Levenshtein similarity with threshold 0.72
    \item $\text{column}(t, q)$ matches query terms against column names
    \item $\text{intent}(t, q)$ applies domain-specific boosts (e.g., boosting sales tables for revenue queries)
    \item $\text{penalties}(t)$ reduces scores for archive tables, empty tables, or system tables
\end{itemize}

The weights $w_i$ were determined empirically through evaluation on a held-out query set. German-language handling includes prefix stripping for common patterns (\texttt{dbo.}, \texttt{vw}, \texttt{tbl}, \texttt{KHK}, \texttt{BS}) to improve matching against natural language queries.

\subsubsection{Performance Impact}

Table~\ref{tab:scout-performance} summarizes the performance improvement achieved through Scout Mode compared to on-demand schema discovery.

\begin{table}[htbp]
    \centering
    \caption{Performance comparison with and without Scout Mode caching.}
    \label{tab:scout-performance}
    \begin{tabular}{lccc}
        \toprule
        \textbf{Operation} & \textbf{Without Scout} & \textbf{With Scout} & \textbf{Improvement} \\
        \midrule
        Schema discovery & 10--50 s & 50 ms & 200--1000$\times$ \\
        Table search & 2--5 s & 10--50 ms & 40--500$\times$ \\
        Full query pipeline & 15--60 s & 0.5--3 s & 20--100$\times$ \\
        \bottomrule
    \end{tabular}
\end{table}

% -----------------------------------------------------------------------------
\subsection{Query Reasoning Pipeline}
\label{subsec:query-pipeline}
% -----------------------------------------------------------------------------

User queries are processed through a multi-stage pipeline that transforms natural language into executable SQL, executes the query, and formats results for presentation. Each stage is designed to be observable and auditable, supporting the traceability principle.

\subsubsection{Pipeline Overview}

Figure~\ref{fig:pipeline} illustrates the query processing pipeline. The pipeline is implemented as a ReAct agent loop, where the language model iteratively reasons about the current state and selects tools to execute until a final answer can be generated.

\begin{figure}[htbp]
    \centering
    \begin{tikzpicture}[node distance=0.5cm]

        % Input
        \node[rectangle, rounded corners=3pt, draw=black!80, fill=gray!15,
              minimum width=4cm, minimum height=0.8cm, font=\small\bfseries]
              (input) {Natural Language Query};

        % Stage 1
        \node[stage, below=0.6cm of input] (s1) {1. Intent Parsing};
        \node[stageout, right=0.3cm of s1.east] {Entities, metrics, filters};

        % Stage 2
        \node[stage, below=0.5cm of s1] (s2) {2. Table Discovery};
        \node[stageout, right=0.3cm of s2.east] {Ranked candidates from Scout};

        % Stage 3
        \node[stage, below=0.5cm of s2] (s3) {3. Schema Retrieval};
        \node[stageout, right=0.3cm of s3.east] {Column names, foreign keys};

        % Stage 4
        \node[stage, below=0.5cm of s3] (s4) {4. SQL Generation};
        \node[stageout, right=0.3cm of s4.east] {Dialect-specific query};

        % Stage 5
        \node[stage, below=0.5cm of s4] (s5) {5. Query Execution};
        \node[stageout, right=0.3cm of s5.east] {Bounded result set};

        % Stage 6
        \node[stage, below=0.5cm of s5] (s6) {6. Response Formatting};
        \node[stageout, right=0.3cm of s6.east] {Natural language answer};

        % Output
        \node[rectangle, rounded corners=3pt, draw=black!80, fill=green!15,
              minimum width=4cm, minimum height=0.8cm, font=\small\bfseries,
              below=0.6cm of s6]
              (output) {Structured Response};

        % Arrows
        \draw[dataarrow] (input) -- (s1);
        \draw[dataarrow] (s1) -- (s2);
        \draw[dataarrow] (s2) -- (s3);
        \draw[dataarrow] (s3) -- (s4);
        \draw[dataarrow] (s4) -- (s5);
        \draw[dataarrow] (s5) -- (s6);
        \draw[dataarrow] (s6) -- (output);

        % Error recovery loop
        \draw[dataarrow, dashed, draw=red!60]
            (s5.west) -- ++(-1.2,0) |- node[pos=0.25, left, font=\footnotesize, text=red!70] {Error Recovery} (s3.west);

    \end{tikzpicture}
    \caption{Query processing pipeline stages. The pipeline is implemented as a ReAct agent loop where stages may be revisited during error recovery (dashed line). Each stage produces observable intermediate outputs for traceability.}
    \label{fig:pipeline}
\end{figure}

\subsubsection{LangGraph Agent Architecture}

The pipeline is implemented using LangGraph, a framework for building stateful, multi-step applications with language models. LangGraph extends LangChain by introducing explicit graph-based control flow, enabling complex agent behaviors while maintaining observability. Figure~\ref{fig:langgraph-architecture} illustrates the agent's internal structure.

\begin{figure}[htbp]
    \centering
    \begin{tikzpicture}[node distance=0.6cm]

        % ReAct Loop Box
        \node[rectangle, rounded corners=5pt, draw=purple!60, fill=purple!5,
              line width=1pt, minimum width=11cm, minimum height=5.5cm] (reactbox) {};
        \node[above=0.1cm of reactbox.north, font=\small\bfseries, text=purple!70] {ReAct Agent Loop (LangGraph)};

        % Think node
        \node[rectangle, rounded corners=3pt, draw=orange!70, fill=orange!15,
              minimum width=2.8cm, minimum height=1.2cm, font=\small,
              align=center] (think) at (-3, 0.5) {\textbf{Think}\\{\footnotesize LLM reasoning}};

        % Act node
        \node[rectangle, rounded corners=3pt, draw=green!60, fill=green!15,
              minimum width=2.8cm, minimum height=1.2cm, font=\small,
              align=center, right=1.5cm of think] (act) {\textbf{Act}\\{\footnotesize Tool execution}};

        % Observe node
        \node[rectangle, rounded corners=3pt, draw=blue!60, fill=blue!15,
              minimum width=2.8cm, minimum height=1.2cm, font=\small,
              align=center, right=1.5cm of act] (observe) {\textbf{Observe}\\{\footnotesize Process result}};

        % State box below
        \node[rectangle, rounded corners=3pt, draw=gray!60, fill=gray!10,
              minimum width=10cm, minimum height=1cm, font=\small,
              below=0.8cm of act] (state) {
              \textbf{MessagesState}: conversation history, tool calls, tool results
        };

        % Arrows in loop
        \draw[dataarrow, line width=1pt] (think) -- node[above, font=\footnotesize] {tool\_call} (act);
        \draw[dataarrow, line width=1pt] (act) -- node[above, font=\footnotesize] {ToolMessage} (observe);
        \draw[dataarrow, line width=1pt, rounded corners=8pt]
            (observe.south) -- ++(0,-0.3) -| node[below, pos=0.25, font=\footnotesize] {next iteration} (think.south);

        % Entry arrow
        \draw[dataarrow, line width=1pt] (-5.5, 0.5) -- (think.west)
            node[pos=0, left, font=\footnotesize] {START};

        % Exit arrow
        \draw[dataarrow, line width=1pt, dashed] (think.north) -- ++(0, 0.8)
            node[right, font=\footnotesize] {no tool\_call $\rightarrow$ END};

        % State arrows
        \draw[dataarrow, draw=gray!50, line width=0.6pt] (think.south) -- (state);
        \draw[dataarrow, draw=gray!50, line width=0.6pt] (act.south) -- (state);
        \draw[dataarrow, draw=gray!50, line width=0.6pt] (observe.south) -- (state);

        % Tools box on right
        \node[rectangle, rounded corners=3pt, draw=green!50, fill=green!5,
              minimum width=2.5cm, minimum height=2.8cm, font=\footnotesize,
              align=left, right=0.5cm of reactbox.east, anchor=west] (tools) {
              \textbf{Bound Tools}\\[3pt]
              \texttt{discover\_tables}\\
              \texttt{get\_schema}\\
              \texttt{get\_column\_index}\\
              \texttt{list\_relations}\\
              \texttt{execute\_query}
        };

        \draw[dataarrow, draw=green!60] (act.east) -- ++(0.5,0) |- (tools.west);

    \end{tikzpicture}
    \caption{LangGraph ReAct agent architecture. The agent iterates through Think-Act-Observe cycles, with state persisted in a \texttt{MessagesState} object. Tools are bound via LangChain's function calling interface and executed during the Act phase.}
    \label{fig:langgraph-architecture}
\end{figure}

\paragraph{ReAct Pattern Implementation.}
The agent implements the ReAct (Reasoning and Acting) pattern~\cite{yao2023react} through LangGraph's \texttt{create\_react\_agent} primitive. Each iteration consists of three phases:

\begin{enumerate}
    \item \textbf{Think}: The language model receives the current conversation state (including prior tool results) and reasons about the next action. It may decide to call a tool or provide a final answer.

    \item \textbf{Act}: If the model requests a tool call, LangGraph executes the corresponding Python function with the specified arguments. Tool definitions use LangChain's \texttt{@tool} decorator, which automatically generates JSON schemas from function signatures and docstrings for the model's function-calling interface.

    \item \textbf{Observe}: Tool outputs are wrapped in \texttt{ToolMessage} objects and appended to the conversation state, providing the model with execution results for the next reasoning step.
\end{enumerate}

\paragraph{State Management.}
LangGraph maintains agent state through a typed state schema that extends the framework's \texttt{MessagesState} base class:

\begin{itemize}
    \item \texttt{messages}: Chronological list of conversation messages (system prompt, user query, assistant responses, tool calls, tool results)
    \item \texttt{question}: Original natural language query
    \item \texttt{tables\_used}: Tables selected during discovery
    \item \texttt{sql\_query}: Generated SQL statement
    \item \texttt{query\_result}: Execution results
    \item \texttt{final\_answer}: Formatted response
\end{itemize}

This state is passed between nodes and persists across iterations, enabling the agent to build upon previous reasoning steps and tool results.

\paragraph{Tool Binding.}
Tools are defined as Python functions with the \texttt{@tool} decorator from LangChain. The decorator extracts type hints and docstrings to generate OpenAI-compatible function schemas:

\begin{verbatim}
@tool
def execute_query(sql: str) -> str:
    """Execute a SQL query against the database.

    Args:
        sql: The SQL query to execute (must be SELECT)

    Returns:
        Query results as formatted markdown table
    """
\end{verbatim}

The language model sees these schemas and can invoke tools by name with structured arguments. LangGraph handles the dispatch, execution, and result routing automatically.

\paragraph{Iteration Control.}
The agent's execution is bounded by a configurable recursion limit (default: 25 iterations) that prevents runaway API costs from infinite loops or excessively complex queries. Each Think-Act-Observe cycle consumes two iterations (one for the agent node, one for the tool node). Typical queries complete in 3--8 iterations; queries requiring error recovery may use 10--15 iterations.

\paragraph{Event Streaming.}
For real-time observability, the agent uses LangGraph's \texttt{astream\_events} API, which emits fine-grained events during execution:

\begin{itemize}
    \item \texttt{on\_chat\_model\_start}: LLM begins processing
    \item \texttt{on\_chat\_model\_end}: LLM response ready (includes any tool calls)
    \item \texttt{on\_tool\_start}: Tool execution begins
    \item \texttt{on\_tool\_end}: Tool execution complete with results
\end{itemize}

These events are captured by the debug logger and can be streamed to the user interface, providing transparency into the agent's decision-making process.

\subsubsection{Stage 1: Intent Parsing}

The first stage extracts structured intent from the natural language query using the language model. The output is a typed structure containing:

\begin{itemize}
    \item \textbf{Primary entities}: The main business objects referenced (e.g., ``products,'' ``customers'')
    \item \textbf{Metrics}: Requested measurements or aggregations (e.g., ``total sales,'' ``count'')
    \item \textbf{Filters}: Constraints on the data (e.g., ``last month,'' ``status = active'')
    \item \textbf{Discovery keywords}: Terms for semantic table search
    \item \textbf{Confidence score}: Model's self-assessed confidence (0.0--1.0)
\end{itemize}

A key architectural insight was that intent should be extracted \emph{once} and reused throughout the pipeline, rather than re-extracted at each stage. Early iterations that performed redundant extraction suffered from keyword inflation, where different extractions produced inconsistent terms that polluted table discovery. The single-extraction approach reduced discovery candidates from 4,715 to under 50, significantly improving downstream accuracy.

\subsubsection{Stage 2: Table Discovery}

Using the discovery keywords from intent parsing, the system queries the Scout catalog to identify relevant tables. The semantic ranking described in Section~\ref{subsec:scout-mode} returns the top-$k$ candidates (default: 5) with relevance scores and explanatory reasoning.

The discovery stage also retrieves foreign key relationships between candidate tables, enabling the subsequent SQL generation stage to construct appropriate joins without requiring the language model to guess relationship paths.

\subsubsection{Stage 3: Schema Retrieval}

Before SQL generation, the exact column names for selected tables are retrieved from the Scout catalog. This prevents a common failure mode where the language model generates plausible but incorrect column names based on general knowledge rather than the actual schema.

\subsubsection{Stage 4: SQL Generation}

With discovered tables, their schemas, and foreign key relationships, the language model generates a SQL query. The system prompt constrains generation to the target database dialect (SQL Server or PostgreSQL) and enforces patterns such as:

\begin{itemize}
    \item Use of \texttt{TOP} (SQL Server) or \texttt{LIMIT} (PostgreSQL) for result bounding
    \item Proper quoting of identifiers containing special characters
    \item Appropriate date functions for temporal filters
    \item Join syntax based on discovered foreign keys
\end{itemize}

\subsubsection{Stage 5: Query Execution}

Generated SQL is validated and executed through the MCP server. Validation includes:

\begin{itemize}
    \item \textbf{Read-only enforcement}: Only \texttt{SELECT} statements are permitted
    \item \textbf{Dangerous keyword rejection}: Statements containing \texttt{DROP}, \texttt{TRUNCATE}, \texttt{DELETE}, or \texttt{EXEC} are rejected
    \item \textbf{Result limiting}: A maximum row count is enforced regardless of query content
    \item \textbf{Timeout enforcement}: Queries exceeding the configured timeout are cancelled
\end{itemize}

\subsubsection{Stage 6: Response Formatting}

Query results are formatted into a natural language response by the language model. The response contextualizes the data relative to the original question and may include summary statistics, notable patterns, or caveats about data interpretation.

\subsubsection{Error Recovery}

The ReAct agent architecture enables automatic error recovery. If SQL execution fails (e.g., due to an invalid column name), the agent observes the error message and can take corrective action---typically retrieving the exact column index and regenerating the query. This recovery loop is bounded by a configurable iteration limit (default: 25 iterations) to prevent runaway API costs.

% -----------------------------------------------------------------------------
\subsection{Production Integration}
\label{subsec:production-integration}
% -----------------------------------------------------------------------------

Deploying the system against production enterprise databases introduces additional architectural considerations around network access, security, and operational constraints.

\subsubsection{Network Architecture}

Enterprise databases are typically isolated within private networks, accessible only through VPN connections. The system accommodates this through a proxy-based architecture where the MCP server runs on a machine with direct database connectivity (potentially within the enterprise network), while the reasoning layer runs externally.

Figure~\ref{fig:network-architecture} illustrates a typical production deployment where the reasoning layer runs on a development workstation (macOS), communicating over HTTPS with an MCP server running on a Windows machine that has VPN access to the production database.

\begin{figure}[htbp]
    \centering
    \begin{tikzpicture}[node distance=1cm]

        % Development Machine
        \node[netnode=blue!70] (dev) {
            \textbf{Development Machine (macOS)}\\[4pt]
            \begin{tikzpicture}
                \node[component=blue!70, minimum width=2cm] (webui) {Web UI};
                \node[component=blue!70, minimum width=2cm, right=0.2cm of webui] (agent) {ReAct Agent};
                \node[component=blue!70, minimum width=2cm, right=0.2cm of agent] (openai) {OpenAI Client};
            \end{tikzpicture}
        };
        \node[right=0.3cm of dev.north east, anchor=north west, font=\footnotesize\ttfamily, text=blue!70] {Port 3000, 5001};

        % Connection label 1
        \node[below=0.4cm of dev, font=\footnotesize] (conn1label) {HTTPS (TLS 1.3) + API Key Authentication};

        % MCP Server
        \node[netnode=green!60, below=1.2cm of dev] (mcp) {
            \textbf{MCP Server (Windows)}\\[4pt]
            \begin{tikzpicture}
                \node[component=green!60, minimum width=1.8cm] (jsonrpc) {JSON-RPC};
                \node[component=green!60, minimum width=1.8cm, right=0.15cm of jsonrpc] (scoutrun) {Scout Runner};
                \node[component=green!60, minimum width=1.8cm, right=0.15cm of scoutrun] (validate) {Query Validator};
                \node[component=green!60, minimum width=1.8cm, right=0.15cm of validate] (vpn) {VPN Client};
            \end{tikzpicture}
        };
        \node[right=0.3cm of mcp.north east, anchor=north west, font=\footnotesize\ttfamily, text=green!60] {Port 8000};

        % Connection label 2
        \node[below=0.4cm of mcp, font=\footnotesize] (conn2label) {VPN Tunnel + ODBC (Read-Only)};

        % Production Database
        \node[netnode=red!50, below=1.2cm of mcp] (db) {
            \textbf{Production Database}\\[4pt]
            \begin{tikzpicture}
                \node[component=red!50, minimum width=2.2cm] (sqlsrv) {SQL Server};
                \node[component=red!50, minimum width=2.2cm, right=0.2cm of sqlsrv] (readonly) {Read-Only User};
                \node[component=red!50, minimum width=2.2cm, right=0.2cm of readonly] (audit) {Audit Logging};
            \end{tikzpicture}
        };
        \node[right=0.3cm of db.north east, anchor=north west, font=\footnotesize\ttfamily, text=red!50] {Port 1433};

        % Arrows
        \draw[dataarrow, line width=1.2pt] (dev) -- (mcp);
        \draw[dataarrow, line width=1.2pt] (mcp) -- (db);

        % Security annotations on left side
        \begin{scope}[on background layer]
            \node[draw=gray!40, dashed, rounded corners=5pt, fit=(dev),
                  inner sep=8pt, label={[font=\footnotesize, text=gray]left:External Network}] {};
            \node[draw=orange!40, dashed, rounded corners=5pt, fit=(mcp),
                  inner sep=8pt, label={[font=\footnotesize, text=gray]left:DMZ / Proxy}] {};
            \node[draw=red!30, dashed, rounded corners=5pt, fit=(db),
                  inner sep=8pt, label={[font=\footnotesize, text=gray]left:Enterprise Network}] {};
        \end{scope}

    \end{tikzpicture}
    \caption{Production deployment architecture showing separation of the reasoning layer (external network) from the data access layer (DMZ/proxy) and enterprise database (internal network). The MCP server acts as a secure gateway enforcing authentication, query validation, and read-only constraints.}
    \label{fig:network-architecture}
\end{figure}

\subsubsection{Security Model}

Security is enforced at multiple architectural layers:

\begin{enumerate}
    \item \textbf{Transport Security.} All communication between the reasoning layer and MCP server uses HTTPS with TLS 1.3. Certificate validation prevents man-in-the-middle attacks.

    \item \textbf{Authentication.} The MCP server requires API key authentication via the \texttt{X-API-Key} header. Keys are stored in environment variables, never in code or configuration files.

    \item \textbf{Query Validation.} The MCP server validates all SQL queries before execution, rejecting any non-SELECT statements or queries containing dangerous keywords.

    \item \textbf{Database Permissions.} The database connection uses a read-only user account with SELECT permissions only. Even if query validation were bypassed, the database would reject modification attempts.

    \item \textbf{Rate Limiting.} The MCP server enforces rate limits (default: 60 requests per minute) to prevent resource exhaustion.

    \item \textbf{Audit Logging.} All queries are logged with timestamps, source identifiers, and execution results for compliance and debugging purposes.
\end{enumerate}

\subsubsection{VPN Latency Accommodation}

VPN connections introduce additional latency (typically 200--500ms round-trip time) that must be accommodated in timeout configurations. The system uses extended timeouts (60 seconds vs. 30 seconds for direct connections) and implements retry logic with exponential backoff for transient network failures.

\subsubsection{Environment Adaptability}

The same application code supports multiple deployment configurations through environment-based configuration:

\begin{itemize}
    \item \textbf{Development}: Direct connection to local PostgreSQL for rapid iteration
    \item \textbf{Staging}: Proxy connection to test SQL Server instance
    \item \textbf{Production}: Proxy connection over VPN to production database
\end{itemize}

This adaptability is achieved through a database client abstraction layer that transparently handles connection mode selection and SQL dialect differences.

% -----------------------------------------------------------------------------
\subsection{Summary}
\label{subsec:architecture-summary}
% -----------------------------------------------------------------------------

The presented architecture addresses the challenges of enabling natural language access to enterprise databases through several key design decisions:

\begin{itemize}
    \item A \textbf{layered architecture} separates concerns between presentation, reasoning, and data access, enabling flexible deployment across network boundaries.

    \item \textbf{Scout Mode} provides autonomous schema modeling, eliminating the need for manual database configuration while enabling sub-second query response times through semantic caching.

    \item A \textbf{ReAct agent pattern} implements the query reasoning pipeline with built-in error recovery, while composite tools reduce the number of reasoning iterations required.

    \item \textbf{Multi-layer security} enforces read-only constraints through query validation, database permissions, and audit logging.

    \item \textbf{Proxy-based production integration} accommodates typical SME network constraints including VPN access requirements.
\end{itemize}

The architectural decisions documented in this section are formalized in the project's Architecture Decision Records (ADRs), which provide additional detail on alternatives considered and rationale for each decision. The complete ADR collection and implementation source code are available in the project repository.

% =============================================================================
% END OF SYSTEM ARCHITECTURE SECTION
% =============================================================================

% -----------------------------------------------------------------------------
% APPENDIX PLACEHOLDERS (to be included in appendix chapter)
% -----------------------------------------------------------------------------
%
% \section{Architecture Decision Records}
% \label{appendix:adrs}
%
% The following ADRs document key architectural decisions:
% \begin{itemize}
%     \item ADR-0005: Model Context Protocol adoption
%     \item ADR-0007: MCP Database Server implementation
%     \item ADR-0014: Scout Mode semantic caching
%     \item ADR-0016: Phase 7+ architecture with Scout and semantic ranking
%     \item ADR-0021: Semantic intent parsing
%     \item ADR-0030: Simple SQL Agent architecture
%     \item ADR-0031: Simple SQL Agent complete architecture (includes Scout integration)
% \end{itemize}
% Full ADR text is available in the project repository.
%
% \section{Scout Catalog JSON Schema}
% \label{appendix:catalog-schema}
%
% [Include JSON schema or representative example]
%
% -----------------------------------------------------------------------------
